% !TeX root = 04_intro_ml_polytech.tex
% !TeX encoding = UTF-8
% !TeX spellcheck = en_GB
\begin{frame}{Neural Networks}{Some history to start\dots}
    Let us start from what we know (classification):
    \begin{equation*}
        \vec{y}
        =
        \only<1>{\sigma}
        \only<2->{\underbrace{\sigma}_{\text{sigmoid}}}
        \left(
        X
        \only<1>{\vec{\beta}}
        \only<2->{\underbrace{\vec{w}}_{\text{weights}}}
        \right) \in \left[ 0, 1 \right]^n
    \end{equation*}

    \vfill \pause

    We can graphically represent this as:

    \begin{center}
        \begin{tikzpicture}
    \node[circle, draw=CEAblack, fill=CEAblack!15] (out) at (2.5, 1.5) {$X \vec{w}$};
    \foreach \x in {0,...,3} {
            \node[circle, draw=CEAblack, fill=CEAblack!15]
            (x\x) at (0, 3-\x) {\footnotesize $x_{\x}$};
            \draw[thin, ->, CEAblack] (x\x) -- (out) node[midway, above, sloped] {\footnotesize $w_{\x}$};
        }
    \node[circle, draw=CEAblack, fill=CEAblack!15] (y) at (5, 1.5) {$\vec{y}$};
    \draw[thin, ->, CEAblack] (out) -- (y) node[midway, above] {\footnotesize $\sigma$};

    % Perceptron label
    \node[anchor=north west, font=\bfseries] at (out.south) {Rosenblatt's perceptron (1958)};
\end{tikzpicture}

    \end{center}
\end{frame}

\begin{frame}{Neural Networks}{Rosenblatt's perceptron briefly\dots}
    Rosenblatt's \textbf{perceptron} (1958) is a simple model of \textbf{artificial neurons}:

    \vfill

    \begin{columns}[c]
        \begin{column}{0.5\textwidth}
            \centering
            {\large \textbf{AND} gate} \\[1em]
            \begin{tikzpicture}[
        ax/.style={->, thick, CEAblack},
        true/.style={circle, thick, draw=CEAgreen, fill=CEAgreen!15, text=CEAgreen},
        false/.style={circle, thick, draw=CEAred, fill=CEAred!15, text=CEAred},
    ]

    \draw[ax] (-2, 0) -- (2, 0) node[below] {$x_1$};
    \draw[ax] (0, -2) -- (0, 2) node[left] {$x_2$};

    \node[true] at (1, 1) {\tiny (T, T)};
    \node[false] at (1, -1) {\tiny (T, F)};
    \node[false] at (-1, 1) {\tiny (F, T)};
    \node[false] at (-1, -1) {\tiny (F, F)};

    % Linear separator
    \draw<2->[thick, CEAblue] (-1, 2) -- (2, -1) node[midway, above, sloped] {\footnotesize $x_1 + x_2 = 1$};
\end{tikzpicture}
 \\
            \only<2->{\textbf{Linear separable} by the perceptron}
        \end{column}
        \hfil
        \begin{column}{0.5\textwidth}
            \centering
            {\large \textbf{XOR} gate} \\[1em]
            
\begin{tikzpicture}[
        ax/.style={->, thick, CEAblack},
        true/.style={circle, thick, draw=CEAgreen, fill=CEAgreen!15, text=CEAgreen},
        false/.style={circle, thick, draw=CEAred, fill=CEAred!15, text=CEAred},
    ]

    \draw[ax] (-2, 0) -- (2, 0) node[below] {$x_1$};
    \draw[ax] (0, -2) -- (0, 2) node[left] {$x_2$};

    \node[false] at (1, 1) {\tiny (T, T)};
    \node[true] at (1, -1) {\tiny (T, F)};
    \node[true] at (-1, 1) {\tiny (F, T)};
    \node[false] at (-1, -1) {\tiny (F, F)};
\end{tikzpicture}
 \\
            \only<2->{\textbf{Not linear separable} by the perceptron}
        \end{column}
    \end{columns}
\end{frame}

\begin{frame}{Neural Networks}{Rosenblatt's perceptron briefly\dots}
    The keyword is \textbf{non-linearity}:
    \begin{block}{Non-linear decision boundaries}
        \centering
        The perceptron must be able to create \textbf{non-linear decision boundaries} to solve complex problems (e.g.\ the XOR gate).
    \end{block}

    \vfill \pause

    How to achieve this? Let us start from $X \in \mathbb{R}^{n \times p}$:
    \begin{equation*}
        \only<2>{X\, W}
        \only<3>{%
            X\, W^{(0)}
        }
        \only<4>{%
            \underbrace{\alert{f^{(0)}}(X\, W^{(0)})}_{h^{(0)}}
        }
        \only<5>{%
        \underbrace{\alert{f^{(1)}}(f^{(0)}(X\, W^{(0)})\, \alert{W^{(1)}})}_{h^{(1)}}
        }
        \only<6>{%
        \underbrace{\alert{f^{(2)}}(f^{(1)}(f^{(0)}(X\, W^{(0)})\, W^{(1)})\, \alert{W^{(2)}})}_{h^{(2)}}
        }
        \only<7>{%
        \vec{y}
        =
        \alert{\sigma}(\dots
        f^{(2)}(f^{(1)}(f^{(0)}(X\, W^{(0)})\, W^{(1)})\, W^{(2)})
        \dots
        \alert{W^{(L)}})
        }
    \end{equation*}
    where
    \only<2>{$W \in \mathbb{R}^{p \times m}$}
    \only<3->{%
    $W^{(0)} \in \mathbb{R}^{p \times m^{(0)}}$
    \only<5->{, $W^{(1)} \in \mathbb{R}^{m^{(0)} \times m^{(1)}}$}
    \only<6->{, $W^{(2)} \in \mathbb{R}^{m^{(1)} \times m^{(2)}}$}
    \only<7->{, \dots, $W^{(L)} \in \mathbb{R}^{m^{(L-1)} \times n}$}
    }

    \only<4->{%
        and $f^{(i)}$ are \textbf{non-linear} functions called \textbf{activation functions} (\emph{activations}).
    }
\end{frame}

\begin{frame}{Neural Networks}{Rosenblatt's perceptron briefly\dots}
    Let us see this graphically:

    \vfill

    \begin{columns}[c]
        \begin{column}{0.6\textwidth}
            \centering
            \begin{tikzpicture}[
        neuron/.style={circle, draw=CEAblack, fill=CEAblack!15},
        activated/.style={circle, draw=CEAred, fill=CEAred!15},
        input/.style={neuron, fill=CEAgreen!15},
        output/.style={neuron, fill=CEAblue!15},
        weight/.style={draw=CEAblack, thin, ->},
        scale=0.8,
        every node/.style={scale=0.8},
    ]

    \foreach \x in {0,...,4} {
            \node[input] (i\x) at (0, 4-\x) {\footnotesize $x_{\x}$};
        }
    \node[output] (o) at (8, 2) {\footnotesize $y$};

    % First layer
    \foreach \x in {0,...,2} {
            \node<2>[neuron] (h1\x) at (1.5, 3-\x) {\tiny $h^{(0)}_{\x}$};
            \node<3->[activated] (h1\x) at (1.5, 3-\x) {\tiny $h^{(0)}_{\x}$};
            \foreach \y in {0,...,4} {
                    \draw<2->[weight] (i\y) -- (h1\x) node[midway, above, sloped] {\only<2>{\tiny $w^{(0)}_{\y,\x}$}};
        }k
    }

    % Second layer
    \foreach \x in {0,...,3} {
            \node<4>[neuron] (h2\x) at (3, 3.5-\x) {\tiny $h^{(1)}_{\x}$};
            \node<5->[activated] (h2\x) at (3, 3.5-\x) {\tiny $h^{(1)}_{\x}$};
            \foreach \y in {0,...,2} {
                    \draw<4->[weight] (h1\y) -- (h2\x) node[midway, above, sloped] {\only<4>{\tiny $w^{(1)}_{\y,\x}$}};
        }
    }

    % Third layer
    \foreach \x in {0,...,2} {
            \node<6>[neuron] (h3\x) at (4.5, 3-\x) {\tiny $h^{(2)}_{\x}$};
            \node<7->[activated] (h3\x) at (4.5, 3-\x) {\tiny $h^{(2)}_{\x}$};
            \foreach \y in {0,...,3} {
                    \draw<6->[weight] (h2\y) -- (h3\x) node[midway, above, sloped] {\only<6>{\tiny $w^{(2)}_{\y,\x}$}};
        }
    }

    % Ellipsis
    \node<7-> (h4) at (6, 2) {\tiny $\dots$};

    % Output layer
    \draw<8->[weight] (h4) -- (o) node[midway, above, sloped] {\only<8>{$\sigma$}};

\end{tikzpicture}

        \end{column}
        \hfil
        \begin{column}{0.4\textwidth}
            \textbf{Input layer:} $X \in \mathbb{R}^{n \times p}$ \\
            \textbf{Output class:} $\vec{y} \in \mathbb{R}^{n}$ \\[1em]

            \onslide<2,4,6>{%
            \textbf{Hidden layer:} $H^{(n)} \in \mathbb{R}^{m^{(n-1)} \times m^{(n)}}$ \\
            }
            \onslide<3,5,7>{%
                \textbf{Activation:} $f^{(n)}$ \\[1em]
            }
            \onslide<8->{%
                \textbf{Final activation:} $\sigma$ \\
            }
        \end{column}
    \end{columns}

    \vfill

    \onslide<9->{%
        A \textbf{neural network} is a \textbf{composition} of linear and \emph{non-linear} activation functions!
        These are often called \textbf{fully connected} (or \emph{dense}) neural networks: each neuron of a layer is connected to all neurons of the next layer.
    }
\end{frame}

\begin{frame}{Neural Networks}{Activation functions in a nutshell\dots}
    Activation functions are non-linear and mostly depend on the \emph{role} of the layer in the network:

    \vfill

    \begin{columns}[T]
        \begin{column}{0.5\textwidth}
            \centering
            {\large \textbf{Hidden layers}} \\[1em]

            \alert<2>{ReLU},
            \alert<3>{Leaky ReLU},
            \alert<4>{Tanh},
            \alert<5>{GELU},
            \alert<6>{SiLU}, etc. \\[1em]

            \begin{tikzpicture}[scale=0.7, every node/.style={scale=0.7}]
    \only<2-3,5->{%
        \begin{axis}[
                axis lines = middle,
                xlabel = {$x$},
                ylabel = {activation},
                xmin = -2, xmax = 2,
                ymin = -0.35, ymax = 2,
                domain = -2:2,
                samples = 100,
                smooth,
                thick,
                grid = major,
                grid style = {dashed, gray!30},
                legend pos = north west
            ]
            \only<2>{%
                \addplot[color=CEAblue, id=ReLU] {max(0, x)};
                \addlegendentry{ReLU(x) = max(0, x)};
            }
            \only<3>{%
                \addplot[color=CEAblue, id=LeakyReLU] {max(0.1*x, x)};
                \addlegendentry{Leaky ReLU(x) = max($\alpha$ x, x)};
            }
            \only<5>{%
                \addplot[color=CEAblue, id=GELU] {0.5*x*(1+tanh(sqrt(2/pi)*(x+0.044715*x^3)))};
                \addlegendentry{GELU(x) = x $\Phi$(x)};
            }
            \only<6>{%
                \addplot[color=CEAblue, id=SiLU] {x/(1+exp(-x))};
                \addlegendentry{SiLU(x) = x $\sigma$(x)};
            }
            \only<7->{%
                \addplot[color=CEAred, id=ReLU] {max(0, x)};
                \addlegendentry{ReLU(x) = max(0, x)};
                \addplot[color=CEAblue, id=LeakyReLU] {max(0.1*x, x)};
                \addlegendentry{Leaky ReLU(x) = max($\alpha$ x, x)};
                \addplot[color=CEAgreen, id=GELU] {0.5*x*(1+tanh(sqrt(2/pi)*(x+0.044715*x^3)))};
                \addlegendentry{GELU(x) = x $\Phi$(x)};
                \addplot[color=CEAyellow, id=SiLU] {x/(1+exp(-x))};
                \addlegendentry{SiLU(x) = x $\sigma$(x)};
            }
        \end{axis}
    }
    \only<4>{%
        \begin{axis}[
                axis lines = middle,
                xlabel = {$x$},
                ylabel = {activation},
                xmin = -5, xmax = 5,
                ymin = -1, ymax = 1,
                domain = -5:5,
                samples = 100,
                smooth,
                thick,
                grid = major,
                grid style = {dashed, gray!30},
                legend pos = north west
            ]
            \addplot[color=CEAblue, id=Tanh] {tanh(x)};
            \addlegendentry{Tanh(x)};
        \end{axis}
    }
\end{tikzpicture}

        \end{column}
        \hfil
        \begin{column}{0.5\textwidth}
            \centering
            {\large \textbf{Output layer}} \\[1em]

            \begin{itemize}
                \item<8-> No activation (regression) $$f(x) = x$$
                \item<9-> Sigmoid (binary classification) $$f(x) = \frac{1}{1 + e^{-x}} \in [0, 1]$$
                \item<10-> Softmax (multi-class classification) $$f(x_i) = \frac{e^{x_i}}{\sum_{j=1}^{K} e^{x_j}} \in [0, 1]$$
            \end{itemize}

        \end{column}
    \end{columns}
\end{frame}

\begin{frame}{Neural Networks}{Are they really black boxes?}
    Neural networks are often called \textbf{black boxes}:

    \vfill

    \begin{center}
        \begin{tikzpicture}[
        neuron/.style={circle, draw=CEAblack, fill=CEAblack!15},
        activated/.style={circle, draw=CEAred, fill=CEAred!15},
        input/.style={neuron, fill=CEAgreen!15},
        output/.style={neuron, fill=CEAblue!15},
        weight/.style={draw=CEAblack, thin, ->},
        scale=0.8,
        every node/.style={scale=0.8},
    ]

    \foreach \x in {0,...,4} {
            \node[input] (i\x) at (0, 4-\x) {\footnotesize $x_{\x}$};
        }
    \node[output] (o) at (8, 2) {\footnotesize $y$};
    \begin{scope}[on background layer]
        \node<2->[set, fit=(i0)(i4), inner sep=2em, label={[set label]Input layer}] {};
        \node<3->[set green, fit=(o), inner sep=2em, label={[set label]Output layer}] {};
    \end{scope}

    % First layer
    \foreach \x in {0,...,2} {
            \node[activated] (h1\x) at (1.5, 3-\x) {\tiny $h^{(0)}_{\x}$};
            \foreach \y in {0,...,4} {
                    \draw[weight] (i\y) -- (h1\x);
                }k
        }

    % Second layer
    \foreach \x in {0,...,3} {
            \node[activated] (h2\x) at (3, 3.5-\x) {\tiny $h^{(1)}_{\x}$};
            \foreach \y in {0,...,2} {
                    \draw[weight] (h1\y) -- (h2\x);
                }
        }

    % Third layer
    \foreach \x in {0,...,2} {
            \node[activated] (h3\x) at (4.5, 3-\x) {\tiny $h^{(2)}_{\x}$};
            \foreach \y in {0,...,3} {
                    \draw[weight] (h2\y) -- (h3\x);
                }
        }

    % Ellipsis
    \node (h4) at (6, 2) {\tiny $\dots$};
    \begin{scope}[on background layer]
        \node<4->[set red, fit=(h10.north west)(h20)(h30)(h12.south west)(h23)(h32)(h4), inner sep=2em, label={[set label]Hidden layers}] {};
    \end{scope}

    % Output layer
    \draw[weight] (h4) -- (o) node[midway, above, sloped] {\only<8>{$\sigma$}};
\end{tikzpicture}

    \end{center}

    \begin{itemize}
        \item<2-> \textbf{Input layer:} this is the layer \emph{you} control (it's the data you feed to the network)
        \item<3-> \textbf{Output layer:} this layer \emph{predicts} the output of the network
        \item<4-> \textbf{Hidden layers:} hidden layers perform \emph{known} operations, but you do not control them
    \end{itemize}
\end{frame}

\begin{frame}{Neural Networks}{How to train a neural network?}
    As for previous models, we define a \textbf{loss function} $\mathcal{L} = \mathcal{L}(\vec{y}, \hat{\vec{y}}; \Theta)$ which\dots
    \begin{itemize}
        \item<2-> \dots depends on the \textbf{true labels} $\vec{y}$ and the \textbf{predictions} $\hat{\vec{y}}$ of the network
        \item<3-> \dots depends on the \textbf{parameters} (weights) $\Theta = \left\{ W^{(1)}, W^{(2)}, \dots, W^{(L)} \right\}$
    \end{itemize}

    \vfill \pause

    We then need to find $$\Theta^* = \arg\min_{\Theta} \mathcal{L}(\vec{y}, \hat{\vec{y}}; \Theta)$$
    which is done by \textbf{gradient descent}:
    \begin{enumerate}
        \item<4-> \textbf{Forward pass:} compute the predictions $\hat{\vec{y}}$ of the network
        \item<5-> \textbf{Backward pass:} compute the gradients of the loss function w.r.t.\ the parameters $\nabla_{W^{(0)}} \mathcal{L}$, $\nabla_{W^{(1)}} \mathcal{L}$, \dots, $\nabla_{W^{(L)}} \mathcal{L}$
        \item<6-> \textbf{Update:} update each of the parameters $\Theta$: $$W^{(\ell)} \gets W^{(\ell)} - \eta \nabla_{W^{(\ell)}} \mathcal{L}$$
    \end{enumerate}
\end{frame}

\begin{frame}{Neural Networks}{Backpropagating errors\dots}
    The training of a neural network is done by \textbf{backpropagation} (1986).
    It uses the fact that
    \begin{center}
        \textbf{neural network} = \textbf{composition of functions}
    \end{center}
    \begin{equation*}
        \hat{y}
        =
        H^{(L)}
        \circ
        H^{(L-1)}
        \circ
        \dots
    \end{equation*}
    where $\circ$ is the composition operator and $H^{(\ell)} = f^{(\ell)}(Z^{(\ell)})$ and $Z^{(\ell)} = H^{(\ell-1)} W^{(\ell)}$.

    \vfill \pause

    We can use the \emph{chain rule}:
    \begin{equation*}
        \footnotesize
        \begin{split}
            \pdv{\mathcal{L}}{W^{(\ell)}_{ij}}
            & =
            \pdv{\mathcal{L}}{\hat{y}}
            \pdv{\hat{y}}{W^{(\ell)}_{ij}} \\
            & =
            \pdv{\mathcal{L}}{\hat{y}}
            \pdv{\hat{y}}{f^{(L)}}
            \pdv{f^{(L)}}{z^{(L-1)}}
            \pdv{z^{(L-1)}}{W^{(\ell)}_{ij}} \\
            & =
            \pdv{\mathcal{L}}{\hat{y}}
            \pdv{\hat{y}}{f^{(L)}}
            \pdv{f^{(L)}}{z^{(L-1)}}
            \pdv{z^{(L-1)}}{f^{(L-2)}}
            \dots
            \pdv{f^{(\ell)}}{z^{(\ell-1)}}
            \pdv{z^{(\ell-1)}}{W^{(\ell)}_{ij}}
            \Rightarrow
            W^{(\ell)}_{ij} \gets W^{(\ell)}_{ij} - \eta \pdv{\mathcal{L}}{W^{(\ell)}_{ij}}
        \end{split}
    \end{equation*}
\end{frame}

\begin{frame}{Neural Networks}{\includegraphics[height=12pt]{01_slides/img/tool.pdf}~Computing backpropagation\dots}
    Consider a \textbf{fully connected} neural network with $2$ hidden layers and ReLU activations trained using a \textbf{MSE} loss on a scalar output $y$:
    \begin{itemize}
        \item \alert<2>{\includegraphics[height=12pt]{01_slides/img/pen-tool.pdf}~Compute the \textbf{derivative of} $\mathcal{L}$ w.r.t.\ the predictions $\hat{y}$}
        \item \alert<3>{\includegraphics[height=12pt]{01_slides/img/brain.pdf}~Does the network need a sigmoid activation function at the output layer?}
        \item \alert<4>{\includegraphics[height=12pt]{01_slides/img/pen-tool.pdf}~Compute the \textbf{derivative of} $\hat{y}$ w.r.t.\ the weights and inputs of the final layer}
        \item \alert<5->{\includegraphics[height=12pt]{01_slides/img/pen-tool.pdf}~Compute the rest of the derivatives and write the update rule}
    \end{itemize}

    \vfill

    \begin{onlyenv}<2>
        \tiny
        \begin{block}{Derivative of the loss function}
            \begin{equation*}
                \pdv{\mathcal{L}}{\hat{y}_i}
                =
                \pdv{\hat{y}_i} \left(\frac{1}{2n} \sum_{i=1}^{n} (y_i - \hat{y}_i)^2\right)
                =
                -\frac{1}{n} (y_i - \hat{y}_i)
            \end{equation*}
        \end{block}
    \end{onlyenv}

    \begin{onlyenv}<4>
        \tiny
        \begin{block}{Derivative of the output layer (scalar!)}
            \begin{equation*}
                \hat{y}_i = \sum\limits_{j} H^{(1)}_{ij} w^{(2)}_{j}
                \implies
                \pdv{\hat{y}_i}{w^{(2)}_{j}} = H^{(1)}_{ij}
                \qand
                \pdv{\hat{y}_i}{H^{(1)}_{jk}} = \delta_{ij} w^{(2)}_k
            \end{equation*}
        \end{block}
    \end{onlyenv}

    \begin{onlyenv}<5>
        \tiny
        \begin{block}{Derivative of the layer $L-1$ (not scalar!)}
            \begin{equation*}
                H^{(1)}_{ij}
                =
                f^{(1)}(\underbrace{\sum\limits_{k} X_{ik} W^{(1)}_{kj}}_{Z^{(1)}_{ij}})
                \implies
                \pdv{H^{(1)}_{ij}}{W^{(1)}_{kl}}
                =
                \pdv{H^{(1)}_{ij}}{f^{(1)}}
                \odot
                \pdv{Z^{(1)}_{ij}}{W^{(1)}_{kl}}
                =
                \theta(Z^{(0)})_{ij}
                \odot
                X_{ik} \delta_{jl}
            \end{equation*}
        \end{block}
    \end{onlyenv}

    \begin{onlyenv}<6>
        \tiny
        \begin{block}{Update rule (I)}
            \begin{equation*}
                w_k^{(2)}
                \gets
                w_k^{(2)}
                -
                \eta \pdv{\mathcal{L}}{w_k^{(2)}}
                =
                w_k^{(2)}
                -
                \eta
                \sum\limits_{i}
                \pdv{\mathcal{L}}{\hat{y}_i}
                \pdv{\hat{y}_i}{w_k^{(2)}}
                =
                w_k^{(2)}
                -
                \frac{\eta}{n}
                \sum\limits_{i}
                (\hat{y}_i - y_i) H^{(1)}_{ik}
                =
                w_k^{(2)}
                -
                \eta\, \mathbb{E}\left[ (\hat{\vec{y}} - \vec{y} ) H^{(1)}\right]
            \end{equation*}
        \end{block}
    \end{onlyenv}

    \begin{onlyenv}<7>
        \tiny
        \begin{block}{Update rule (II)}
            \begin{equation*}
                \begin{split}
                    W_{kl}^{(1)}
                    \gets
                    W_{kl}^{(1)}
                    -
                    \eta \pdv{\mathcal{L}}{W_{kl}^{(1)}}
                    & =
                    W_{kl}^{(1)}
                    -
                    \eta
                    \sum\limits_{i,m,n}
                    \pdv{\mathcal{L}}{\hat{y}_i}
                    \pdv{\hat{y}_i}{H^{(1)}_{mn}}
                    \pdv{H^{(1)}_{mn}}{W^{(1)}_{kl}}
                    =
                    W_{kl}^{(1)}
                    -
                    \frac{\eta}{n}
                    \sum\limits_{i,m,n}
                    (\hat{y}_i - y_i) \delta_{im} w_n^{(2)}
                    \theta(Z^{(0)})_{mn} \odot X_{mk} \delta_{nl}
                    \\
                    & =
                    W_{kl}^{(1)}
                    -
                    \frac{\eta}{n}
                    \sum\limits_{i}
                    (\hat{y}_i - y_i) w_l^{(2)} \theta(Z^{(0)})_{il} \odot X_{ik}
                    =
                    W_{kl}^{(1)}
                    -
                    \eta \mathbb{E}\left[ (\hat{\vec{y}} - \vec{y} ) (\theta(Z^{(0)}) \odot X w^{(2), \top} ) \right]
                \end{split}
            \end{equation*}
        \end{block}
    \end{onlyenv}

    {\tiny Call $\theta(x)$ (1 if $x > 0$, 0 otherwise) the derivative of the ReLU function.}
\end{frame}

\begin{frame}{Backpropagation}{What are the advantages?}
    As we just saw, the \textbf{update} depends only on standard building blocks
    \begin{equation*}
        \pdv{Z^{(\ell)}}{W^{(\ell)}}, \quad
        \pdv{H^{(\ell)}}{f^{(\ell)}}, \quad \text{etc.}
    \end{equation*}

    \vfill

    These can be \textbf{saved} during the forward pass, and simply \textbf{used} in the backward pass:
    \begin{itemize}
        \item Compute the forward pass (i.e.\ compute $\hat{y}$)
        \item Store \emph{layer-by-layer} the blocks containing the derivatives of the layer
        \item Compute the backward pass by \textbf{multiplying} blocks from the loss to the layer
    \end{itemize}
\end{frame}
